\section{Native verification of the rotated trace}\label{sec:native-details}

The following proofs retain the controller, cell-visit, and exception-cost
obligations used in \cref{sec:cd}.

\subsection{Physical row rotations}\label{sec:details-cd-rotated-native}

\begin{informalproof}[Proof of \cref{prop:cd-rotated-native}]
Use the two banks and the paid setup and page executor of
\cref{prop:ne-native-shear}.  Add two reusable physical row-buffer
tapes, denoted $P_{\mathrm{buf}},S_{\mathrm{buf}}$, each with an origin and a movable
end marker.  Their symbols have fixed bit, origin, and end tracks.

We first specify a paid split-and-concatenate primitive.  Given a public
cut $a\in[0,W]$, clear the previous cut marker by a complete outward and
return sweep of the row clock.  Initialize an origin-based binary
countdown to $a$, move the clock one cell after each decrement, mark the
cut when the countdown becomes zero, and rewind the clock.  The countdown
does not perform a full zero test at each decrement: mark its highest
one, update that mark when a borrow removes it, and remember zero when
no one remains.  The total digit traversal in a countdown of $a$ is
$O(a+\log(a+1))$, including returns to the low digit.  Initial parameter
arithmetic uses the complete-sweep compiler of the preceding section.
The cut therefore costs $O(W+\operatorname{polylog}L)$, and all these
head motions are home-based.

To split a sequential $W$-bit source row, scan its bits with the row clock.
In one transition per bit copy into $P_{\mathrm{buf}}$ before the cut and into
$S_{\mathrm{buf}}$ at or after it, advancing only the source, active buffer,
and clock.  A fixed mode changes at the cut flag.  Mark the two buffer
ends and rewind the buffers and clock completely.  An empty buffer has
its end marker at zero.  On reuse, writing a bit clears any former end
flag at that cell; if a shorter string is stored, its new first end marker
stops subsequent reading.  Cells beyond that end are neither read nor
erased.  To concatenate in rotated order, scan $S_{\mathrm{buf}}$ and then
$P_{\mathrm{buf}}$, copying to the sequential destination or directly appending
to the real output.  Each bit transfer is one native transition, and the
buffers are fully rewound afterward.  There is no long-lived payload
register and at most one output bit per transition.  The entire primitive
costs $O(W+\operatorname{polylog}L)$.

For input row $i$, split at $W-ib$ and concatenate $S_{\mathrm{buf}}$ followed
by $P_{\mathrm{buf}}$ into its bank row.  This is a right rotation by $ib$
bits, so raw column $c$ reaches internal column $c+i\pmod d$ without
changing its within-entry bit offset.  Process all rows in order, loading
$A$ and then $B$, with no output appends.  Each bank head moves monotonically
through its bank and pauses during the buffer operations.  Write its
row-start flags on the first concatenated bit of each row and its end
marker on completing the bank; no full-bank setup sweep is used.  After loading,
rewind both banks completely.  Execute the $J$ shear pages exactly as in
\cref{prop:ne-native-shear}, including suffix scans and complete
synchronized rewinds.

Finally process bank rows $i=0,\ldots,d-1$ in order.  Split the bank row
at $ib$, and append $S_{\mathrm{buf}}$ followed by $P_{\mathrm{buf}}$.  This is a
left rotation by $ib$ bits and implements $R^{-1}$.  A source bank row is
copied into buffers before its output is emitted.  In particular, even
the unchanged $A$ output is delayed until after every core page.  We do
not give its syntactic cone an unjustified restriction to $A$ inputs.
The resulting $B$ output row is, explicitly,
\[
 B_j[c,s]+\sum_{i:A_iB_j\in E(H_m)} A_i[c+j-i\pmod d,s],
\]
with addition modulo two.

There are $d$ input and $d$ output primitives.  All parameter calculations,
cut marking, countdown borrows and returns, data copies, buffer rewinds,
bank rewinds, and appends are included in
\[
 O\bigl(L+dW+d\operatorname{polylog}L
              +m\operatorname{polylog}L+\operatorname{polylog}L\bigr)
 =O(L).
\]
The fixed page family, compiler, counters, buffers, and clocks use a fixed
number of tapes and fixed control.  End flags and public row indices
determine all movements, independently of the payload.  Banks traverse
monotonically from origin to end, with pauses, and return completely;
buffers and clocks make origin-to-end excursions; counters return to their
own origins.  Thus the entire physical trace is home-based.  During
concatenation a bank may pause inside a row, so the row-start parking
property of the unwrapped executor is not claimed here.

It remains to certify the actual paths.  Fix one edge and internal lane.
Its raw source bit is first written to an input-buffer cell, then copied
from that same cell to $A_i[gb+s]$.  The core XOR copies its contribution
to $B_j[gb+s]$; in the final wrapper that bank cell is copied to an
output-buffer cell and then read for the actual append.  Follow all
consecutive visits in these four residence intervals: input buffer to
bank copy, $A$ bank to XOR, $B$ bank to output-buffer copy, and output
buffer to append.  Include every parked rewrite.  Consecutive visits
are previous-visit edges, and successive intervals meet at the same copy
or XOR transition.  Reversing their concatenation gives a literal DAG
path with the displayed physical endpoints.

At any time each active bank carrier participates in at most $\Delta$
chosen paths.  The same holds for a live buffer carrier, because it
belongs to one source or one final target bit.  Reuse of a buffer cell
does not accumulate congestion: its residence intervals for different
rows are disjoint, and each path has left the cell before it is reused.
The input and output phases are disjoint, and buffering an $A$ output
carries none of the selected paths.  A step scans one cell of each of the
four payload tapes, so the union bound gives congestion $4\Delta$,
including steps at which a path changes carrier and all parked endpoints.
\end{informalproof}

\subsection{Extracting all row words from the native trace}\label{sec:details-cd-row-words}

\begin{informalproof}[Proof of \cref{lem:cd-row-words}]
Call everything through parameter, endpoint-list, and clock initialization
setup.  Thereafter the successive macro phases are
\[
 I,\ R_0,\ F_1,\ R_1,\ldots,\ F_J,\ R_J,\ O.
\]
Here $I$ is the complete input wrapper, $O$ the complete output wrapper,
$R_0$ the post-load synchronized rewind, $F_a$ the forward part of page
$a$, and $R_a$ its synchronized rewind.  There are $M_{\mathrm{ph}}$ phases,
each occurring once.  A backward edge crossing phases strictly decreases
the phase index.

In a wrapper phase, row episode $i$ contains its parameter calculation,
cut marking, countdown, split, concatenation, tests, and returns.  Episodes
occur in increasing global row order and do not interleave.  In a forward
page, a scan event consists of a whole $W$-bit $A$ skip, $B$ skip, or
matched $A/B$ scan and its following clock return.  Number events in
chronological order, ignoring intervening controller ticks.  There are
at most $2m=d$ events in one page.  Each row of either bank is active in
exactly one event and its active events are in increasing order.  Let
\[
 \tau_{a,A},\tau_{a,B}:[m]\longrightarrow[E_a]
\]
be the corresponding strictly increasing injections.  Their values agree
at a matched scan.  One common event index is used even when one bank is
parked; the parked cursor is not an independently selectable row.

We define a sparse set $X$ of marked real time vertices.  First mark the
first post-setup state vertex and the first post-setup visit to each work
cell touched in setup.  The setup inventory, not its time, is
$O(W+d\operatorname{polylog}L)$.  These marks block every backward edge
entering setup.  Next mark the first and last ticks of every macro phase
and wrapper row episode, all synchronized-rewind turn and end-test ticks,
all slow controller, parser, list, and arithmetic ticks in core pages,
and the first and last ticks of each core row scan and clock return.
The core slow work totals $O(d\operatorname{polylog}L)$.  Wrapper countdowns
are not included in this wholesale marking: their time can be linear.

On each bank mark moving visits to row-boundary cells, and the first tick
of each maximal constant-position run of at least two ticks.  There are
$O(d)$ such runs per fixed phase: a bank pauses at the constant number of
operation boundaries in each wrapper row, and during core row or controller
operations.  A long countdown gives one bank wait, not a new wait at each
decrement.  This rule also covers an interior bank pause between the two
concatenation segments.

For each macro phase, mark its first visit to each cell of the row clock
and the two row buffers.  They occupy only $[0,W]$, giving $O(M_{\mathrm{ph}}W)$
marks.  Finally, for each wrapper row episode mark its first visit to
each cell of every other controller or work tape.  Those tapes use only
polylogarithmically many distinct cells in that episode: metadata heads
park at their origins during wrappers, and the scratch ruler, arithmetic,
and countdown use their bounded polylogarithmic windows.  This costs
$O(d\operatorname{polylog}L)$.  We have not marked first visits to every
bank cell in every phase, nor every clock cell in every row.  Either
would incorrectly charge a linear number of vertices.  The stated rules
give
\[
 |X|\le C[W+d(1+\log_2(L+2))^h].
\]

Assign each unmarked post-setup vertex a chart: its global row-episode
index in $I$ or $O$, its event index in $F_a$, and the common local bank
row $\lfloor\operatorname{position}/W\rfloor$ in $R_a$.  At good rewind
vertices both banks move in lockstep through interior cells.  We check
that every same-phase edge between good vertices changes chart by at
most one.

For state edges this follows from chronological row or event order.
For a wrapper controller tape, an edge into an earlier episode starts at
the first visit to that cell in the current episode and was marked.
All remaining such edges stay in their row, including every countdown
borrow and return.  The row clock receives a complete outward-and-return
sweep while clearing the cut in every wrapper row.  A clock predecessor
therefore lies in the same or preceding row episode.  Its first macro-phase
touch is already marked.

For buffers, every used prefix, including its end, is visited in the row
episode.  During $I$, the $P_{\mathrm{buf}}$ lengths $W-ib$ decrease and the
$S_{\mathrm{buf}}$ lengths $ib$ increase; during $O$ these monotonicities are
reversed.  If a cell has an earlier visit within a wrapper phase, this
monotonicity ensures that the immediately preceding episode also visited
it, unless it has already been visited in the current episode.  Position
zero is visited even for an empty buffer.  The first touch in each macro
phase was marked, so stale earlier-phase visits cause no exception.

For a wrapper bank, any stationary visit after the first has predecessor
$t-1$, while the first tick of a nontrivial parking run was marked.  An
isolated moving interior visit belongs to its physical row episode.
The bank progresses monotonically through that wrapper, so an earlier
same-phase visit to that cell is in the same row.  Moving row-boundary
visits were marked.  These facts cover the possible interior pauses.

At good core forward vertices, controller tapes and buffers are stationary;
after the marked first fast tick their predecessors are $t-1$.  The same
holds for a parked bank after its marked first parking tick.  Each active
interior bank cell is visited once in the forward page, so any other
predecessor leaves the phase.  The moving clock visits each interior cell
on the scan and return of every event, giving a predecessor in the same
or preceding event.  In a synchronized rewind, each bank cell is visited
once, all other heads are stationary, and state edges change local row by
at most one.  Turn and end-test vertices were marked.  This completes the
same-phase edge audit.

Now consider an edge with good endpoints in different phases.  State
crossings are blocked by phase marks, clock and buffer crossings by their
first macro-phase touches, and controller crossings by per-row first
touches in wrappers or by stationary/slow-tick rules in the core.  Setup
crossings have also been blocked.  Only bank cell edges remain.  The
current endpoint cannot be the first tick of a parking run or a moving
row boundary.  It therefore has the chart of its actual active interior
bank row.  At the earlier endpoint, any wrapper pause at that interior
cell belongs to the same physical row; a page or rewind has its unique
active visit.  Thus both endpoints refer to one and the same local bank
row $u\in[m]$.

For bank $H\in\{A,B\}$ define the strict injections
\[
 \begin{array}{lll}
 \iota_{I,A}(u)=\iota_{O,A}(u)=u,&
 \iota_{I,B}(u)=\iota_{O,B}(u)=m+u,\\
 \iota_{R_a,A}(u)=\iota_{R_a,B}(u)=u,&
 \iota_{F_a,A}(u)=\tau_{a,A}(u),&
 \iota_{F_a,B}(u)=\tau_{a,B}(u).
 \end{array}
\]
A backward bank edge from phase $F$ to an earlier phase $F'$ acts exactly
by $\iota_{F',H}\circ \iota_{F,H}^{-1}$ on its chart, with the inverse interpreted
on its image.  This factorization uses a common physical cell and does
not assume that an increasing partial map preserves distances.

Every append occurs in $O$ with chart equal to its physical output row.
After input loading the raw head reads blank.  During $I$ it reads a bit
of the current row while splitting, or waits at the first bit of the next
row during concatenation and controller work.  Exclude the $d$ raw
row-start coordinates.  Every other raw attachment at a good vertex then
has raw row equal to its chart.

Exclude also the demand labels of outputs in the $D$-step forward
neighborhood of $X$ in the full DAG.  There are at most $B_D|X|$ of them.
A free cone of any remaining output can meet neither $X$ nor setup,
because its paths have length at most $D$.  The physical output exclusions
are converted through $\Pi$ to coordinate labels before being united with
the raw row-start exclusions.  This gives the set $Z$ in the statement.

Break a surviving output-to-raw path at its cross-phase edges.  Each
same-phase segment has chart displacement at most $D$.  Each cross-phase
edge contributes the two strict partial injections just displayed.  Since
phase indices strictly decrease, there are at most $M_{\mathrm{ph}}-1$
crossings.  Enlarging every local displacement and every zero gap between
these injections to a band of radius $D+2$ gives the asserted words.
There are at most $2\cdot3^{M_{\mathrm{ph}}-2}$ phase-and-bank choices for paths
from $O$ to $I$; the larger value $K=3^{M_{\mathrm{ph}}}+1$ is convenient and
also covers empty or padded choices.  The bounds on $K,\ell_{\max}$ are independent
of depth.  Finally $W/L=1/d$ and
$d\operatorname{polylog}L/L\to0$, proving the fixed-$D$ error estimate.
\end{informalproof}
